% Simulation de désintégrations aléatoires comparée à la loi de décroissance (L'essentiel, chapitre 6)
\begin{tikzpicture}
\begin{axis}[otMesures,otAxesBout,width=100mm,height=52mm,xmin=0,xmax=42,ymin=0,ymax=1100,
  xtick={0,10,...,40},ytick={0,250,...,1000},xlabel={intervalle $k$},ylabel={nombre de noyaux restants},
  legend style={at={(0.97,0.97)},anchor=north east,draw=otRule,fill=white,font=\otLight\fontsize{7.5}{9}\selectfont,cells={anchor=west}}]
\addplot[otGreenDark,line width=0.8pt] table[x index=0,y index=2] {figures/donnees/06-simulation.dat};
\addlegendentry{loi théorique $N_0\,(1-p)^k$}
\addplot[only marks,mark=*,mark size=1.3pt,otNavy] table[x index=0,y index=1] {figures/donnees/06-simulation.dat};
\addlegendentry{simulation (1\,000 noyaux)}
\end{axis}
\end{tikzpicture}
\par\smallskip\legende{Les points simulés fluctuent autour de la loi théorique : la désintégration est aléatoire, mais un grand nombre de noyaux suit une loi régulière.}
